% Lösungen Einheit 2 von 4 – Wahrscheinlichkeit einstufig (zusammenbau v0.1)
\einheitenkopf{Einheit 2 von 4 · Wahrscheinlichkeit einstufig}

\erg{11}{a) eine Stufe}
\erg{12}{a) $\frac{3}{5}$}
%% TODO Lösung der Erklärzeile 13g) fehlt in der Bank
\erg{13}{a) möglich: 35 Lose (30 Nieten und 5 Gewinne), günstig: 5 Gewinne \quad b) $\frac{1}{6}$ \quad c) $\frac{1}{8}$ \quad d) $\frac{1}{4}$ \quad e) $\frac{1}{6}$ \quad f) $\frac{1}{6}$ \quad g) \ldots \quad h) $\frac{2}{8} = \frac{1}{4}$ \quad i) $\frac{3}{9} = \frac{1}{3}$ \quad j) $\frac{2}{6} = \frac{1}{3}$ \quad k) $\frac{9}{50} = 0{,}18 = 18\,\%$ \quad l) $\frac{15}{45 + 15} = \frac{15}{60} = \frac{1}{4}$ \quad m) $\frac{6}{8} = \frac{3}{4}$ \quad n) $0{,}75 \cdot 4 = 3$ Felder \quad o) 3 rote, 2 blaue, 1 gelbes Feld \quad p) Nein: $\frac{3}{16}$, denn es sind nur noch 16 Muffins da, davon 3 mit Chili. \quad q) $\frac{93}{600} = 0{,}155$; $\frac{1}{6} \approx 0{,}167$ – die Werte liegen nahe beieinander. \quad r) Nein: Beide Augenzahlen haben die Wahrscheinlichkeit $\frac{1}{6}$. \quad s) Ja: $\frac{4}{50}$, denn Lose mit Endziffer 5 sind 205, 215, …, 695, also 50 Lose; davon enden auf 45: 245, 345, 445, 545, 645 – ohne 445 bleiben 4.}

13o)

\kreisdiagramm{rot/1,rot/1,rot/1,blau/1,blau/1,gelb/1}

\erg{14}{a) mittlere Dose}
\erg{15}{a) $\frac{1}{6} \cdot 300 = 50$-mal}
\erg{16}{a) Fehler: Die leeren Batterien gehören auch zu allen Batterien. Richtig: $\frac{3}{50}$. \quad b) Die Felder sind nicht gleich groß; Rot nimmt den halben Kreis ein, also ist die Wahrscheinlichkeit $\frac{1}{2}$. \quad c) $\frac{35}{100} = \frac{7}{20} = 0{,}35$ \quad d) $\frac{1}{8} \cdot 40 = 5$ Figuren erwartet – sicher ist das nicht, es können mehr oder weniger sein.}
